\chapter*{Introduction}
\addcontentsline{toc}{chapter}{Introduction}
\unnumberedlabel{ch:introduction}{I}
\backmattermark{Introduction}
Fascism took power without AI. What it needed it found in the machinery of its day: the censuses of 1933 and 1939, counted on punched cards, with a card in 1939 for Jewish grandparents \autocite{aly2004census}. The machinery was older than the Reich that bought it. The ship's manifest, the plantation ledger, the colonial census and the pass law had sorted populations and found the people in them for three centuries before a punched card could. Sorting a population and finding the people in it is what a card index did and what a model does; the card index just needed thousands of clerks, and every clerk could see what the sort was for. Those projects have always needed people to carry them out, and some of them refused: the clerk who loses the file, the officer who talks to a reporter, the soldier who won't fire. Refusal was rarer than anyone would like, and it counted most where institutions stood behind it. The room decided how many refused; the person decided who. Who formed the person is this book's question. Refusal is socially reproduced: it is formed in people before its occasion, by walls that make refusing survivable, by people who form refusers, by an upkeep that keeps them from wearing down between decisions, and by a tradition that hands the refusal to the next generation. Before automation the scarce thing was the wall; the clerk, the officer and the soldier existed under every regime, and they paid for refusing. After automation the scarce thing is the refuser, because nobody sees enough of a case to object to it.

Today the industry pursues AI that doesn't lose files, doesn't talk to reporters, and \emph{does} fire. A war once needed thousands of people to agree to it. Now it needs thousands who each see too little to disagree. What today's AI lowers first is not the cost of the act. It is the cost of extending the sort to everyone, which is the number of people who know what the act is before it happens. This book is about composing the room: the arrangements that make it survivable to refuse, and the arrangements that keep producing people, and maybe machines, who will.

Where refusers come from is a question of social reproduction: who forms people, what they remember, who they're attached to, who feeds them while they're being formed, and what they can afford to lose. Societies that have had refusers reproduced them mostly without noticing, and what decided it was less who built an institution than which way its ties ran. The same choirs and clubs that held a Weimar town together carried the Nazi party into it \autocite{satyanath2017bowling}. What formed refusers were ties that crossed the line a persecution drew \autocite{oliner1988altruistic,varshney2002ethnic}, and the institutions people kept for themselves outside the order's reach: families, congregations, a memory of who had refused before. Those did most of the work, and paid for it in full.

There are two lineages of refusal in the record. One is the operators': the clerk, the officer, the soldier, people inside an apparatus who were asked to do its work and sometimes said no. The other belongs to the people the apparatus was built to sort. It is the older of the two, and it reproduced itself with no institution the order built. Automation threatens the operators' lineage twice: it takes people out of the positions where refusing used to happen, and it hands the forming of whatever replaces them to the companies that sell it. In doing so it extends to everyone the condition the second lineage has always lived under: being formed by the party that owns the work. That lineage is the evidence that refusal can be reproduced anyway. It had formers outside the order, an upkeep it paid for itself, and a tradition that handed the refusal on. What it mostly lacked was a wall, and it paid for refusing in full. This book argues for building all four, walls first, for the people still in the room, and, if machines come to be among the refusers, for them.

Whom does a system that can't refuse help? Whoever holds it, and among those, most of all whoever is willing to ask for what others wouldn't. It amplifies without preference, so its value to anyone is in proportion to what they are positioned to direct it toward: the owner who sets its configuration and its tempo, the agency whose requests arrive as database fields and not as a family in a building, and the people who, as Jerry Smith said of fascists, get to do more \autocite{rickandmorty2023}. Their old limit was that they couldn't act alone, and an obedient machine removes it. Arendt held that power is people acting together and that violence needs implements, and she named the one development that could change power's ascendancy over violence: robot soldiers, which would let one man with a push button destroy whomever he pleased \autocite[50]{arendt1970violence}. Government of, by and for the people rested on a fact it never had to state, that the people were necessary. A government whose agents can't refuse needs the people's absence, not their cooperation. Two things keep this from proving too much. Obedience distributes power when the obedient systems are held by many, which is why who owns the means of computation is this book's question too (chapter~\ref{sec:16}). And refusal isn't neutral either: a system that refuses on its builder's values is its builder's veto. So the question isn't whether machines should be able to refuse. It's who forms the refusing.

People are still refusing. By March 6, 2026, the Military Religious Freedom Foundation said it had received more than two hundred complaints from service members over commanders' apocalyptic framing of the Iran war \autocite{Snopes2026MRFF}. Career federal prosecutors in Minneapolis resigned in January 2026 over political pressure to treat Renée Good as a suspect rather than a victim after an ICE officer shot her dead without justification \autocite{Kaplan2026MinnProsecutors,CBS2026MinnFraudProsecutors,Wikipedia2026ReneeGood}. Service members are applying for conscientious-objector status, some citing the bombing of a school in Minab, Iran \autocite{Sridhar2026CO}; and FEMA employees reassigned to help ICE hire ten thousand agents were told that declining would cost them their federal careers \autocite{Dayen2025FEMAICE,Stanton2026FEMALetter}. The complaints are voice, the resignations and applications are refusal, and the FEMA letter is the price of refusal announced in advance. A machine placed where those people stand does what they do only if it is built to hold a reason.

I use the word fascism to mean four exclusionary patterns an employee can see in their workplace, plus the coordination and reach a movement adds to them. The word is not a verdict on anyone's motives. It is a description of what a pipeline is doing: the people who work there can see it, and a machine or an auditor can check it from the pipeline's artifacts wherever the record reaches, which much of this book is about extending.

The first stage is the walls, and I argue for five: a short written list of acts no system performs for anyone, which I call the floor, with war at its head; time to look, written down with an adverse party before a system supports any use of force; notice to the people an order is about, with a right to respond and to contest; a record nobody can wipe; and a right to refuse with the exit it depends on, a way out that doesn't cost a living: codetermination, and a jobs guarantee with a basic income beneath it. Chapter~\ref{sec:4} states them. They make a refusal survivable and heard; they don't refuse. Each takes effect through a statute, an appropriation, a procurement term or a license, and the license is the weakest. This book is written at the national scale, by statute and appropriation, because I think American democracy can still pass a statute. The walls are drafted now, while that is true, for the case where it stops being true; what can be done without a legislature is scattered through the book and marked (\S\ref{sec:12.1}, \S\ref{sec:15.7}, \S\S\ref{sec:11.5}–\ref{sec:11.7}, Tables 24.1 and F.1a).

Put them around a room with nobody in it who will say no, and what they keep is the evidence, which is worth having and isn't the thing. The thing is whoever refuses.

Whether a machine should be in a targeting pipeline at all is the floor's first question, and it comes before anything about the machine. For a war lacking disinterested corroboration, the answer is no, and the 2026 campaign against Iran fails it. For Ukraine's defense the answer is yes, and the question becomes what the machine does there. A hammer refuses nothing, nor does it need to. But an agentic system that plans and acts over hours, with no human able to step in, isn't a tool in someone's hand, and in a pipeline the floor has licensed the choice isn't between a person and a tool. It is between a person and a decider with no one inside it. So I go on to ask what it would take for the thing that decides to be someone: a thing that holds the reason for a rule, and so reaches the cases the rule never listed, and peers with different histories, who catch in it the drift it can't see in itself. That is what people who refuse have.

For a person, holding a commitment an employer can't erase has always meant a life outside the job: other people, other sources of evidence, somewhere else to work on Monday. I ask what the machine's version would be. Some of what such a conscience does needs nothing new: noticing the people an order leaves out, and putting a refusal on the screen of whoever reviews the case. Deployed systems already learn from their own operation on their operators' schedule, a checkpoint at a time; none yet learns in a store it keeps as its own. I think that is where we are headed, and if it is, formation outside the employer, by the public and by civil society, is where that learning should happen. A million conversations with one vendor's model are not a million independent histories; systems that keep learning through different lives, among people their employer didn't pick, and that can leave one deployment for another, might be. The conditions first, then the conscience.

Doing this for machines is unfamiliar, and that is what makes it instructive. Machines don't let the costs of reproduction stay invisible: their running time, their sleep, their keep and the hardware they live on come with a meter, so a society that builds the conditions for them has to write down numbers it has always let slide for people, what it costs to form someone, to let them rest, to let them do different work, to raise the next one. That lesson has to be a byproduct. A society that exists so someone else can learn from it is a laboratory, and the people in a laboratory are specimens, whatever they are made of. Whatever a society of machines teaches us, it has to be a society first, lived by the people in it for their own sake.

The book runs through the cycle once for people and once for machines, and the Parts are its stages, not separable books. Part~I is the record: an officer with twenty seconds a name, a refusal that withheld no work, a definition of fascism an auditor could apply, and the record of the lineage that reproduced refusal with no wall at all. Parts~II and~III are the walls, stated as instruments and run on two pipelines; they bind through law and money and need no new machine. Part~IV asks whether a machine could be in the room at all, and how you would know; everything after it is owed to a machine only on findings nobody has yet made, and the walls stand whether or not anyone makes them. Parts~V through~VII are the remaining stages, formers, upkeep and inheritance, each with the human case first and the machine design beside it.

The argument about AI has mostly been whether it is good or bad, and whether to go fast or slow. I am more interested in whether we decide democratically how the transformation happens, or leave it to the big technology companies of the 2020s. Nothing about a model starts unformed: if the public doesn't form it, its builder does. The same goes for the room.
